\section{Что осталось гипотетическим и как продолжать}

Установленные здесь утверждения решают конкретные задачи для конкретных
моделей и классов дифференцирований. Следующая таблица фиксирует задачи,
которые нельзя записывать в список решённых.

\begin{longtable}{L{0.30\textwidth}L{0.64\textwidth}}
\toprule Задача & Точный статус и следующий проверяемый шаг \\ \midrule
\endhead
Все гармонические вложения & Даны два бесконечных семейства и конечные
проверки. Нужен общий переход между классами вычетов либо новое
конструктивное выпрямление. Ограниченный перебор не является классификацией.\\[6pt]
Комплексное выпрямление $H(3,7,8)$ и $H(4,5,7)$ & Не построено.
Можно искать локально конечное поле в найденных аффинных пространствах
кандидатов. Отсутствие такого поля в этих степенях, если будет доказано,
не исключит поля большей степени.\\[6pt]
Полный ML новых ядер & Пока имеется дихотомия $k$ или $k[t]$.
Для результата $k$ достаточно одного LND, двигающего $t$.
Для результата $k[t]$ надо доказать, что каждый LND фиксирует $t$;
ограниченный поиск этого не доказывает.\\[6pt]
Однородный $E$ со срезом $t$ & Условия $E(t)=1$ и сохранения ядра
ещё не дают LND. Обнуление итераций на образующих является точным
положительным сертификатом; неудача конечного числа итераций не
доказывает отсутствия подходящего поля.\\[6pt]
Формальные координаты & Построены все конечные утолщения и пополнение.
Не доказано, что координаты обрываются в полиномы: это отдельная
глобальная задача.\\[6pt]
Question 4.21 Dasgupta--Gaifullin & Формулы фильтраций получены,
но минимальность трёхшагового биррационального представления не
установлена. Нужно исключить все допустимые более короткие представления.\\[6pt]
Минимальные степени выпрямителей & Степени предъявленных автоморфизмов
являются верхними оценками. Минимумы степеней срезов не заменяют эту задачу.\\[6pt]
Вещественный спуск & Вещественное выпрямление рассматриваемых нетривиальных
узлов невозможно. Полиномиальность соответствующего вещественного
ядра без сохранения выделенного $t$ этим сама по себе не исключается.\\[6pt]
Общие задачи сокращения и линеаризации & Отвергнут один предложенный
кандидат; общие задачи и все возможные действия не исследованы полностью.\\
\bottomrule
\end{longtable}

\subsection*{Почему поиск LND со срезом мотивирован}
В кольце $B_\mu$ элемент $t$ выделен самим вложением, а специальный слой
$t=0$ уже имеет полиномиальное описание. Если удастся предъявить
однородное дифференцирование $E$ степени $-1$ с $E(t)=1$ и доказать его
локальную нильпотентность, формула среза даст
\[
 \rho_E(f)=\sum_{j\ge0}\frac{(-t)^j}{j!}E^j(f),\qquad
 B_\mu=(\ker E)[t].
\]
Сумма конечна именно благодаря LND. При однородности эта конструкция
возвращает координаты исходного вложения, как доказано в \S11.
Поэтому дополнительное требование LND отвечает точной геометрической
цели. Оно не является технической деталью, автоматически следующей
из линейных уравнений на коэффициенты.

Для конечного набора образующих $b_1,\ldots,b_m$ достаточно предъявить
целые $N_i$ и тождества $E^{N_i}(b_i)=0$. Формула Лейбница тогда даёт
локальную нильпотентность на всей алгебре. Для поиска нормального
локально конечного поля альтернативный положительный сертификат ---
конечномерное устойчивое пространство, содержащее образующие, из \S12.
Оба способа дают проверяемое доказательство в случае успеха,
но здесь не доказана общая граница времени или степени для их поиска.

\subsection*{Роль групп Кремоны}
Биррациональные замены вида
\[
 (x,y,z)\longmapsto (x,y,m(x)z-q(x,y))
\]
дают эффективный язык переноса полей и координат после локализации.
Возвращение к полиномиальному кольцу требует точной делимости и проверки
центра модификации. Установленный запрет для естественного центра
семейства $H(3,3n+1,3n+2)$ показывает, что формального рационального
обратного отображения недостаточно. Практическая программа состоит
в поиске других центров, в вычислении образных идеалов и в сравнении
различных факторизаций. Общая классификация подгрупп Кремоны здесь
не заявляется.

\subsection*{Связь с теоремой Абъянкара--Моха}
Прямое применение плоской теоремы возможно после предъявления
координатной поверхности, содержащей кривую: на этой поверхности
вложенная прямая становится плоским вложением, к которому применяется
\cite{AM,Suzuki}. Найти такую поверхность в пространстве ---
содержательная недостающая часть. Эквивалентные критерии через
координатный параметр, градуированное LND и нормальное локально конечное
поле позволяют атаковать именно эту часть разными методами.
Не следует подменять её уже известной плоской теоремой.

\subsection*{Порядок чтения}
Для введения в LND полезно начать с \emph{First Principles} и
\emph{Further Properties of LNDs} книги Freudenburg \cite{Freudenburg},
затем перейти к главам \emph{Dimension Three}, \emph{Algorithms},
\emph{Makar-Limanov and Derksen Invariants} и
\emph{Slices, Embeddings and Cancellation}. Для настоящего проекта
следующий круг чтения --- \cite{vdER} (точные гипотезы и кольцо ядра),
\cite{Masuda} (эквивариантная структура), \cite{EVM} (относительные
координаты), \cite{DG} (итерации LND и сложность),
\cite{FF,KP} (гармонические кривые и узлы).

\section{Воспроизводимость и состав приложения}

Все вычислительные файлы помещены в каталог \file{computations/}.
Комментарии в поставляемом коде даны на английском языке. Исходные
версии сохранены в архиве происхождения; отдельно проверяется, что
редактирование комментариев и документационных строк Python не
меняет абстрактное синтаксическое дерево исполняемой части.
Это полезный контроль упаковки, но не доказательство правильности
алгоритма.

\begin{longtable}{L{0.29\textwidth}L{0.65\textwidth}}
\toprule Каталог & Назначение \\ \midrule
\endhead
\file{baseline/} & Исходные выпрямители, точные проверки взаимно
обратных отображений, напечатанных компонент и исходных сертификатов.\\[5pt]
\file{legacy_appendix/} & Вычисления опубликованного аппендикса
от 26 октября, включая локальные зависимости для автономного запуска.\\[5pt]
\file{continuation/} & Четырнадцать проверок продолжения: минимумы
срезов, степени кратности, формальные координаты, ограничения
переноса, фильтрации LND и относительный ML.\\[5pt]
\file{euler/} & Нормальные поля, ранги линейных систем, двойственные
сертификаты несовместности и осевые препятствия для выбранных полей.\\[5pt]
\file{audit_checks/} & Проверка отождествления кандидата из
Conjectures 5.1--5.3 с исходной моделью и исходным срезом.\\[5pt]
\file{verification_logs/} & Стандартный вывод, ошибки и время
каждой программы повторного прогона.\\[5pt]
\file{provenance/} & Исходные версии, контрольные суммы и сведения
о происхождении материалов.\\
\bottomrule
\end{longtable}

Точка входа --- \file{computations/run_all_checks.py}.
Запуск следует выполнять из каталога \file{computations/}:
\begin{quote}\ttfamily\small
python -m pip install -r requirements.txt\\
python run\_all\_checks.py
\end{quote}
Свежий запуск от 28 октября 2026 года завершился успешно: \textbf{35 из 35
программ}. Все основные алгебраические проверки выполняются точно, без замены
тождеств численным приближением. Итоговый машинный протокол
\file{verification_summary.json} фиксирует результаты и команды.
Число программ не равно числу теорем: одна программа может проверять
много тождеств, а бесконечные семейства требуют приведённых выше
математических доказательств.

Код и исторические протоколы GAP/Macaulay2 включены для независимой
проверки, но эти системы в текущем окружении не запускались.
Их старые успешные протоколы явно отделены от нового Python-прогона.
Длительные поисковые эксперименты сохраняются как исходники и данные;
они не подменяют краткие проверяющие программы.

\subsection*{Совместимость с репозиторием}
Базовая версия проверена по коммиту
\file{b954f3609e900761fddc9c8ff03d8301e9604cab}.
Все 75 файлов опубликованного аппендикса совпали с Git-объектами этой
версии. Обновление добавляет только отдельный каталог
\file{updates/2026-10-28/} и указатель на него внутри прежнего
аппендикса. Основной препринт и прежние сертификаты не заменяются.

Содержательно уточняются две прежние исследовательские перспективы.
Во-первых, относительный ML старого аппендикса предлагался условно;
\S10 показывает, что ожидаемая нетривиальность при данном определении
невозможна. Во-вторых, новая граница степени $\ge5$ для $H(4,5,7)$
относится к координатам, сохраняющим одну из исходных переменных.
Общий поиск абсолютной координаты со степени $4$ остаётся обоснованным.
Такое уточнение не требует изменения старого архивного текста,
но должно сопровождать его дальнейшее цитирование.

\subsection*{Проверка читателем}
Сначала следует проверить формулировку конкретной модели и поле
коэффициентов; затем запустить сертификаты выпрямителя; после этого
проверить конструктивный переход к ядру и срезу. Для новых общих
утверждений основным объектом проверки остаётся математическое
доказательство. При обнаружении ошибки необходимо указать модель,
утверждение, команду запуска и полный протокол. Автор приглашает
читателей к такой проверке и к развитию исследовательской программы.
